\documentclass[10pt,a4paper]{amsart}
\usepackage[margin=19mm]{geometry}
\usepackage{amsmath,amssymb,mathtools}
\usepackage[scr=rsfs,cal=euler]{mathalfa}
\usepackage{xfrac}
\usepackage[unicode,hidelinks,bookmarks=false]{hyperref}
\newcommand{\ie}{i.e.\ }
\newcommand{\cf}{cf.\ }
\newcommand{\resp}{respectively\ }
\newcommand{\Subsection}{\subsection}
\providecommand{\nobreakdash}{\nobreak\hbox{-}}
\setlength{\emergencystretch}{2em}
\multlinegap0pt
\usepackage[shortlabels]{enumitem}
\usepackage{booktabs}
\usepackage{amssymb}
\usepackage{nicefrac}
\usepackage{algorithm}\usepackage{algpseudocode}
\usepackage{multirow}
\usepackage{xcolor}
\usepackage{tikz}
\usepackage{natbib}
\newtheorem{assumption}{Assumption}
\newtheorem{theorem}{Theorem}
\newtheorem{lemma}{Lemma}
\newcommand{\E}{\mathbb{E}}
\newcommand{\R}{\mathbb{R}}
\newcommand{\ba}{\mathbf{a}}
\newcommand{\bb}{\mathbf{b}}
\newcommand{\bp}{\mathbf{p}}
\newcommand{\bw}{\mathbf{w}}
\newcommand{\cD}{\mathcal{D}}
\newcommand{\cP}{\mathcal{P}}
\title{A Two-Layer Framework for Joint Online Configuration Selection and Admission Control}
\author{Owen Shen, Haoran Xu, Yinyu Ye, Peter Glynn, Patrick Jaillet}
\date{Source: arXiv:2602.07663v1 (CC BY 4.0)}
\begin{document}
\maketitle

\noindent\emph{Source and licence.} A Two-Layer Framework for Joint Online Configuration Selection and Admission Control. Version arXiv:2602.07663v1 (CC BY 4.0), distributed by arXiv under the Creative Commons Attribution 4.0 International licence, http://creativecommons.org/licenses/by/4.0/, as recorded in the arXiv licence metadata for this exact version and shown on the abstract page https://arxiv.org/abs/2602.07663v1. Full licence notice: \texttt{licenses/arxiv-2602.07663-CC-BY-4.0.txt}.\par\smallskip

\noindent\textit{Editorial scope.} Counted unit: the article's theorem thm:pd\_equilibrium (source0.tex lines 251--264). The article's complete original proof of it is delivered with it (source0.tex lines 1041--1128, one \texttt{proof} environment of 88 lines, which the article places away from the statement). No mathematics is written, repaired or re-derived: the statement and the proof are the article's own lines, in the article's own notation, and no retained line is substituted or re-flowed.  Retained uncounted as the source-local context the proof needs: lines 174--177; lines 178--184; lines 186--189; lines 191--194; lines 216--222; lines 238--248; lines 893--899; lines 954--962; lines 976--987; lines 1030--1033.  Nothing was omitted from any retained span, so the package carries no omission record.  No retained line contains a citation, so the wrapper carries no bibliography and no bibliography entry.  No retained line contains a figure environment, an image-inclusion command or drawing code, so no image is delivered and the package declares no asset dependency.  Editorial formatting only, in addition: an AMS \texttt{amsart} preamble that restates the article's own declarations and macros used by the retained lines, namely the environments \texttt{assumption}, \texttt{theorem}, \texttt{lemma} on separate counters, together with the standard packages those lines need.  The retained environments print in this document as Assumption 1, Theorem 1, Lemma 1 in order of appearance, whereas the article prints the same items with the numbering of its own full document.  The complete native source of the article as distributed in the arXiv e-print for this exact version is bundled unchanged as \texttt{sources/194196/source0.tex}.
\begin{assumption}[Stationarity and Independence]
\label{assump:iid}
If an algorithm select configuration $\theta$ at time period $t$, $(r_t,\ba_t)$ is drawn independently from $\cD_\theta$. In addition, the distribution $\cD_\theta$ is time-invariant for all $\theta\in\Theta$.
\end{assumption}

\begin{assumption}[Boundedness]
\label{assump:bounded}
There exist constants $R_{\max}, A_{\max} > 0$ such that for all $\theta$ and $(r, \ba) \sim \cD_\theta$:
\[
0 \le r \le R_{\max}, \quad \|\ba\|_\infty \le A_{\max} \quad \text{a.s.}
\]
\end{assumption}

\begin{assumption}[Budget Scaling]
\label{assump:budget-scaling}
There exists $\bb\in \mathbb{R}_+^d$ such that $B=T\cdot \bb$.
\end{assumption}

\begin{assumption}[Continuous Valuation]
\label{assump:contivalue}
For any $\theta\in\Theta$, if $(r,\ba)\sim \cD_{\theta}$, then the conditional distribution of $r$ given $\ba=\alpha$ is continuous for any $\alpha$ in the support of $\ba$.
\end{assumption}

\begin{equation}
\label{eq:mix-primal}
\begin{split}
V^{\mathrm{mix}}(\bb) := \max_{\bw \in \Delta_K} \max_{\{x_\theta\}} \Big\{ & \sum_{\theta \in \Theta} w_\theta \E_\theta[r\,x_\theta(r, \ba)] \\
& \Big|\; \sum_{\theta \in \Theta} w_\theta \E_\theta[\ba\,x_\theta(r, \ba)] \le \bb \Big\}.
\end{split}
\end{equation}

\begin{theorem}[Primal--Dual Form of the Mixed Fluid Oracle]
\label{thm:mix-dual}
Under Assumptions~\ref{assump:bounded} and~\ref{assump:budget-scaling}, let $b_{min}$ be the smallest component of $\bb$ (in particular, $b_{\min} > 0$) and $P_{max}$ be $2R_{max}/b_{min}$. With price domain $\cP = [0, P_{\max}]^d$,
\begin{equation}
\label{eq:mix-dual}
\begin{split}
V^{\mathrm{mix}}(\bb) &= \max_{\bw \in \Delta_K} \min_{\bp \in \cP} \left\{ \langle \bp, \bb \rangle + \sum_{\theta} w_\theta g_\theta(\bp) \right\} \\
&= \min_{\bp \in \cP} \left\{ \langle \bp, \bb \rangle + \max_{\theta \in \Theta} g_\theta(\bp) \right\}.
\end{split}
\end{equation}
\end{theorem}

\begin{assumption}[No-Tie Condition]
\label{assump:no-ties}
For all $\theta \in \Theta$ and all $\bp \in \cP$,
\[
\Pr_{(r,\ba) \sim \cD_\theta}\big(r = \langle \bp, \ba \rangle\big) = 0.
\]
\end{assumption}

\begin{lemma}[Upper Box Constraint is Inactive at Minimizers]
\label{lem:upper-box-inactive}
Assume $0 \le r \le R_{\max}$ a.s.\ and $b_{\min} > 0$. Fix any $\bw \in \Delta_K$ and consider the convex function
\[
F_{\bw}(\bp) := \langle \bp, \bb \rangle + \sum_{\theta} w_\theta g_\theta(\bp), \qquad \bp \in \R_+^d.
\]
Then every minimizer of $\min_{\bp \ge 0} F_{\bw}(\bp)$ satisfies $\|\bp\|_\infty \le R_{\max}/b_{\min}$.
In particular, if $P_{\max} > R_{\max}/b_{\min}$, then every minimizer of $\min_{\bp \in \cP} F_{\bw}(\bp)$ lies in the strict interior of the upper box constraints, i.e., $p_i < P_{\max}$ for all $i$.
\end{lemma}

\begin{lemma}[Differentiability and Gradient Formula for $g_\theta$]
\label{lem:grad-g-pop}
Assume Assumption~\ref{assump:bounded} and Assumption~\ref{assump:no-ties} (no ties).
Then for each $\theta$, $g_\theta$ is convex and continuously differentiable on $\cP$, and
\[
\nabla g_\theta(\bp) = -h_\theta(\bp) = -\E_\theta\!\left[\ba \, \mathbf{1}\{r > \langle \bp, \ba \rangle\}\right].
\]
Consequently, for any $\bw \in \Delta_K$,
\[
\nabla_{\bp} L(\bw, \bp) = \bb - H(\bw, \bp).
\]
\end{lemma}

\begin{lemma}[Nonnegative Dot Product Complementarity is Coordinatewise]
\label{lem:coordwise-complementarity}
If $u, v \in \R_+^d$ and $\langle u, v \rangle = 0$, then $u_i v_i = 0$ for every $i \in [d]$.
\end{lemma}

\begin{theorem}[Primal--dual optimality (saddle/KKT conditions)]
\label{thm:pd_equilibrium}
Assume Assumptions~\ref{assump:iid}--\ref{assump:contivalue}.
A pair $(\bw^\star,\bp^\star)\in\Delta_K\times\cP$ is a saddle point of $\max_{\bw \in \Delta_K} \min_{\bp \in \cP}L(\bw,\bp)$
(and hence attains $V^{\mathrm{mix}}(\bb)$) if and only if:
\begin{enumerate}[label=(\roman*),leftmargin=*,noitemsep]
\item \textbf{(Envelope support)} $\mathrm{supp}(\bw^\star)\subseteq A(\bp^\star)$.
\item \textbf{(Feasibility)} $H(\bw^\star,\bp^\star)\le \bb$ (componentwise).
\item \textbf{(Complementarity)} $\langle \bp^\star,\bb-H(\bw^\star,\bp^\star)\rangle=0$.
\end{enumerate}
Moreover, the mixture $\bw^\star$ together with threshold admission $x_{\theta}(r,\ba)=\mathbf{1}\{r>\langle \bp^\star,\ba\rangle\}$ for all $\theta\in\Theta$,
 is primal-optimal for~\eqref{eq:mix-primal} and
$V^{\mathrm{mix}}(\bb)=L(\bw^\star,\bp^\star)$.
\end{theorem}

\begin{proof}[Proof of Theorem~\ref{thm:pd_equilibrium}]
Assume Assumptions~\ref{assump:bounded}--\ref{assump:no-ties}, and assume $P_{\max} > R_{\max}/b_{\min}$ so that by Lemma~\ref{lem:upper-box-inactive} the upper box constraints are inactive at any minimizer of $\bp \mapsto L(\bw, \bp)$ for any fixed $\bw$.

We prove the equivalence between:
\begin{itemize}
\item[(S)] $(\bw^\star, \bp^\star)$ is a saddle point of $L$, i.e.,
\[
L(\bw, \bp^\star) \le L(\bw^\star, \bp^\star) \le L(\bw^\star, \bp) \quad \forall \bw \in \Delta_K, \forall \bp \in \cP;
\]
\item[(KKT)] conditions (i)--(iii) in the theorem statement.
\end{itemize}

\textbf{(S) $\Rightarrow$ (KKT).}

\emph{Step 1 (Envelope support).}
Fix $\bp^\star$. The map $\bw \mapsto L(\bw, \bp^\star) = \langle \bp^\star, \bb \rangle + \sum_{\theta} w_\theta g_\theta(\bp^\star)$ is linear in $\bw$ over the simplex. Therefore, any maximizer $\bw^\star \in \arg\max_{\bw \in \Delta_K} L(\bw, \bp^\star)$ must place all its mass on indices attaining the maximum coefficient $g_\theta(\bp^\star)$:
\[
\mathrm{supp}(\bw^\star) \subseteq A(\bp^\star) := \arg\max_{\theta} g_\theta(\bp^\star).
\]
This is condition (i).

\emph{Step 2 (First-order optimality for the $\bp$-minimization).}
Because $(\bw^\star, \bp^\star)$ is a saddle, $\bp^\star \in \arg\min_{\bp \in \cP} L(\bw^\star, \bp)$.
By Lemma~\ref{lem:upper-box-inactive}, $\bp^\star$ lies in the interior of the upper box constraints, so the effective constraint set is only $\bp \in \R_+^d$.

By Lemma~\ref{lem:grad-g-pop}, $L(\bw^\star, \bp)$ is convex and differentiable in $\bp$ with $\nabla_\bp L(\bw^\star, \bp) = \bb - H(\bw^\star, \bp)$. The KKT condition for minimizing a convex differentiable function over $\R_+^d$ is
\[
\mathbf{0} \in \nabla_\bp L(\bw^\star, \bp^\star) + N_{\R_+^d}(\bp^\star),
\]
where $N_{\R_+^d}(\bp^\star)$ is the normal cone to $\R_+^d$ at $\bp^\star$.
Equivalently, there exists $v \in N_{\R_+^d}(\bp^\star)$ such that
\[
\mathbf{0} = \bb - H(\bw^\star, \bp^\star) + v.
\]
The normal cone satisfies: $v_i = 0$ if $p_i^\star > 0$, and $v_i \le 0$ if $p_i^\star = 0$.
Thus:
\[
p_i^\star > 0 \Rightarrow b_i - H_i(\bw^\star, \bp^\star) = 0, \qquad p_i^\star = 0 \Rightarrow b_i - H_i(\bw^\star, \bp^\star) \ge 0.
\]
Hence $H(\bw^\star, \bp^\star) \le \bb$ componentwise, which is condition (ii). Also, $b_i - H_i(\bw^\star, \bp^\star) \ge 0$ and $p_i^\star \ge 0$ imply $\langle \bp^\star, \bb - H(\bw^\star, \bp^\star) \rangle \ge 0$. Moreover, because the coordinatewise implications above include $p_i^\star(b_i - H_i(\bw^\star, \bp^\star)) = 0$ for each $i$, we obtain
\[
\langle \bp^\star, \bb - H(\bw^\star, \bp^\star) \rangle = 0,
\]
which is condition (iii).

\textbf{(KKT) $\Rightarrow$ (S).}

Assume (i)--(iii).

\emph{Step 1 ($\bw^\star$ is a best response to $\bp^\star$).}
For fixed $\bp^\star$, $\bw \mapsto L(\bw, \bp^\star)$ is linear over $\Delta_K$, and its maximum value equals $\langle \bp^\star, \bb \rangle + \max_{\theta} g_\theta(\bp^\star)$.
Condition (i) implies $\sum_\theta w_\theta^\star g_\theta(\bp^\star) = \max_{\theta} g_\theta(\bp^\star)$, hence for all $\bw \in \Delta_K$,
\[
L(\bw, \bp^\star) \le L(\bw^\star, \bp^\star).
\]

\emph{Step 2 ($\bp^\star$ is a best response to $\bw^\star$).}
By Lemma~\ref{lem:grad-g-pop}, $L(\bw^\star, \bp)$ is convex differentiable in $\bp$.
Since (ii) gives $\bb - H(\bw^\star, \bp^\star) \in \R_+^d$ and (iii) gives $\langle \bp^\star, \bb - H(\bw^\star, \bp^\star) \rangle = 0$, Lemma~\ref{lem:coordwise-complementarity} implies $p_i^\star(b_i - H_i(\bw^\star, \bp^\star)) = 0$ for each $i$.
Equivalently:
\[
p_i^\star > 0 \Rightarrow b_i - H_i(\bw^\star, \bp^\star) = 0, \qquad p_i^\star = 0 \Rightarrow b_i - H_i(\bw^\star, \bp^\star) \ge 0.
\]
This is exactly the KKT condition $\mathbf{0} \in \nabla_\bp L(\bw^\star, \bp^\star) + N_{\R_+^d}(\bp^\star)$ for minimizing $L(\bw^\star, \cdot)$ over $\R_+^d$ (upper box inactive by our choice of $P_{\max}$).
Therefore $\bp^\star \in \arg\min_{\bp \in \cP} L(\bw^\star, \bp)$, and for all $\bp \in \cP$,
\[
L(\bw^\star, \bp^\star) \le L(\bw^\star, \bp).
\]

Combining Steps 1 and 2 yields the saddle inequalities, so $(\bw^\star, \bp^\star)$ is a saddle point.

\textbf{Primal optimality of threshold admission and value identity.}
Let $x(r, \ba) := \mathbf{1}\{r > \langle \bp^\star, \ba \rangle\}$ and define the achieved (fluid) reward
\[
U(\bw^\star, \bp^\star) := \sum_\theta w_\theta^\star \E_\theta\!\left[r \mathbf{1}\{r > \langle \bp^\star, \ba \rangle\}\right].
\]
Using the identity $r \mathbf{1}\{r > \langle \bp^\star, \ba \rangle\} = \langle \bp^\star, \ba \rangle \mathbf{1}\{r > \langle \bp^\star, \ba \rangle\} + (r - \langle \bp^\star, \ba \rangle)_+$, we obtain
\[
U(\bw^\star, \bp^\star) = \langle \bp^\star, H(\bw^\star, \bp^\star) \rangle + \sum_\theta w_\theta^\star g_\theta(\bp^\star).
\]
Hence
\[
L(\bw^\star, \bp^\star) = \langle \bp^\star, \bb \rangle + \sum_\theta w_\theta^\star g_\theta(\bp^\star) = U(\bw^\star, \bp^\star) + \langle \bp^\star, \bb - H(\bw^\star, \bp^\star) \rangle.
\]
By condition (iii), the last inner product is zero, so $U(\bw^\star, \bp^\star) = L(\bw^\star, \bp^\star)$.
Condition (ii) implies the threshold rule is feasible for the primal constraint.
Finally, since $(\bw^\star, \bp^\star)$ is a saddle point, $L(\bw^\star, \bp^\star)$ equals the minimax value $V^{\mathrm{mix}}(\bb)$ by Theorem~\ref{thm:mix-dual}. Therefore the mixture $\bw^\star$ together with threshold admission attains $V^{\mathrm{mix}}(\bb)$ and is primal-optimal for~\eqref{eq:mix-primal}.
\end{proof}

\end{document}
